\documentclass[11pt]{article}
\usepackage{preamble}
\setlength{\parskip}{4pt}

\begin{document}

\daybanner{AP Statistics \textperiodcentered\ Unit 3 \textperiodcentered\ Topic 3.13 \textperiodcentered\ B: 2026-12-17 \textperiodcentered\ E: 2027-01-29 \textperiodcentered\ TEACHER KEY}{One Result, Two Claims}

\sectionbanner{CED TETHER}

{\small
\begin{itemize}[leftmargin=1.5em,itemsep=2pt,topsep=2pt]
  \item \textbf{Skill 3.E} --- Calculate a test statistic and find a p-value, provided conditions for inference are met.
  \item \textbf{Skill 4.B} --- Interpret statistical calculations and findings to assign meaning or assess a claim.
  \item \textbf{Skill 4.E} --- Justify a claim using a decision based on significance tests.
  \item \textbf{EU VAR-6} --- The normal distribution may be used to model variation.
  \item \textbf{LO VAR-6.K} --- Calculate an appropriate test statistic for the difference of two population proportions. [Skill 3.E]
  \item \textbf{EK VAR-6.K.1} --- The test statistic for a difference in proportions is: z = [($\widehat p_1$ - $\widehat p_2$) - 0] / [$\surd$($\widehat p_c$(1-$\widehat p_c$)) $\cdot$ $\surd$(1/$n_1$ + 1/$n_2$)], where $\widehat p_c$ = ($n_1$$\widehat p_1$ + $n_2$$\widehat p_2$) / ($n_1$ + $n_2$).
  \item CLARIFYING STATEMENT: The formulas for test statistics do not appear explicitly on the AP Statistics Formula Sheet. However, these formulas do not need to be memorized, as they can be constructed based on the general test statistic formula and the standard error formulas provided on the formula sheet.
  \item \textbf{EU DAT-3} --- Significance testing allows us to make decisions about hypotheses within a particular context.
  \item \textbf{LO DAT-3.C} --- Interpret the p-value of a significance test for a difference of population proportions. [Skill 4.B]
  \item \textbf{EK DAT-3.C.1} --- An interpretation of the p-value of a significance test for a difference of two population proportions should recognize that the p-value is computed by assuming that the null hypothesis is true, i.e., by assuming that the true population proportions are equal to each other.
  \item \textbf{LO DAT-3.D} --- Justify a claim about the population based on the results of a significance test for a difference of population proportions. [Skill 4.E]
  \item \textbf{EK DAT-3.D.1} --- A formal decision explicitly compares the p-value to the significance $\alpha$. If the p-value $\leq$ $\alpha$, then reject the null hypothesis, $H_0$: $p_1$ = $p_2$, or $H_0$: $p_1$ - $p_2$ = 0. If the p-value > $\alpha$, then fail to reject the null hypothesis.
  \item \textbf{EK DAT-3.D.2} --- The results of a significance test for a difference of two population proportions can serve as the statistical reasoning to support the answer to a research question about the two populations that were sampled.
\end{itemize}
}
\textbf{Essential question:} What does a significant test against zero establish, and what practical claim remains untested?\par
\textbf{DOK-3 skill:} 4.E \textperiodcentered\ \textbf{FRQ pattern:} {\small\texttt{zero-\allowbreak{}test-\allowbreak{}vs-\allowbreak{}practical-\allowbreak{}threshold}}\par
\textbf{DOK rationale:} Part (c) coordinates a two-sided test decision, effect-size comparison, and the claim actually tested to adjudicate competing reports, specify the missing analysis, and prevent a practically overstated conclusion.\par

\frameworkphaseheader{Do Now (first take) $\rightarrow$ Explore (video + follow-along) $\rightarrow$ Exit (finish + turn in)}{1 $\rightarrow$ 3}{45}{%
\begin{itemize}[leftmargin=1.5em,itemsep=2pt,topsep=2pt]
  \item Board slide up at the bell. Let students compare the observed gap and sample sizes before calculating.
  \item Begin the 6.11 video follow-along at minute 5.
  \item Return to the board slide at minute 33. Circulate and require separate statements about the zero-difference test and five-point deployment standard.
  \item Collect at the door. Score part (c) only, E/P/I.
\end{itemize}
}{%
\begin{itemize}[leftmargin=1.5em,itemsep=2pt,topsep=2pt]
  \item Choose Rowan's or Salma's conclusion and cite one number.
  \item Complete the follow-along about pooled standard error, test statistic, p-value, and conclusion.
  \item Finish (a)--(c), revise the first take in the exit line, and turn in the sheet.
\end{itemize}
}{%
\begin{itemize}[leftmargin=1.5em,itemsep=2pt,topsep=2pt]
  \item Which null difference did this test actually use?
  \item What does 0.041 describe under that null model?
  \item How does the observed two-point gap compare with the five-point deployment standard?
  \item What additional interval or test would address the five-point claim directly?
\end{itemize}
}{Solo teacher. Circulate ~5 pairs. Require the test decision and a separate practical-threshold judgment.}

\begin{callout}{\IconWarn\ WATCH FOR}{calloutred}
\begin{itemize}[leftmargin=1.5em,itemsep=2pt,topsep=2pt]
  \item Using unpooled standard errors in a null test of equal proportions.
  \item Interpreting the p-value as the probability that the population proportions are equal.
  \item Concluding that rejection of zero proves an improvement of at least five percentage points.
\end{itemize}
\end{callout}

\sectionbanner{SCORING PART (c) --- E / P / I}

\textbf{Expected elements}
\begin{itemize}[leftmargin=1.5em,itemsep=2pt,topsep=2pt]
  \item \textbf{uses-formal-decision} (required) --- Uses $p\approx0.041<0.05$ to reject equality
  \item \textbf{compares-effect-to-threshold} (required) --- Compares the observed 0.020 gap with the 0.050 standard
  \item \textbf{distinguishes-claims} (required) --- Explains that a zero test does not establish a difference above 0.05
  \item \textbf{states-reader-consequence} (required) --- Explains that the rejected report confuses detectability with practical sufficiency
  \item \textbf{specifies-needed-evidence} (required) --- Requires a 95 percent one-sided lower bound above 0.05
\end{itemize}

{\small\renewcommand{\arraystretch}{1.25}
\noindent\begin{tabular}{|p{0.5in}|p{5.9in}|}
\hline
\textbf{E} & Uses the test decision and both effect-size numbers, diagnoses the overstatement, and specifies the 95 percent lower-bound criterion. \\ \hline
\textbf{P} & Distinguishes significance from practical importance but omits a number, consequence, or interval criterion. \\ \hline
\textbf{I} & Treats $p<0.05$ as proof of a five-point effect or confuses the null value with the threshold. \\ \hline
\end{tabular}}

\clearpage
\focusbanner{TODAY'S PROBLEM --- annotated}

Suppose a utility has 500,000 online accounts. It randomly selects 10,000, assigns 5,000 to message A and 5,000 to B, and records on-time payment. The pre-data question asks whether the population proportions differ at $\alpha=0.05$; deployment requires A to improve payment by at least 5 percentage points.\par\smallskip \textbf{Rowan:} ``I rejected equality, so A has proved an important improvement of at least 5 points. Deploy A.''\par \textbf{Salma:} ``Rejecting equality supports a nonzero difference, but does not automatically establish the separate 5-point requirement.''\par
\begin{center}
\textbf{On-time payment}\par\smallskip
\begin{tabular}{|l|c|c|c|}
\hline
\textbf{Message} & \textbf{n} & \textbf{Paid} & \textbf{Other} \\ \hline
\textbf{A} & 5,000 & 3,050 & 1,950 \\ \hline
\textbf{B} & 5,000 & 2,950 & 2,050 \\ \hline
\end{tabular}
\end{center}
\begin{firsttakebox}{FIRST TAKE (5 min)}
Whose report is more defensible, Rowan's or Salma's? Cite the observed gap or the claim tested.\par
\answer{Not graded. The large samples make a two-point difference statistically detectable; students must decide what the specified deployment rule still requires.}
\end{firsttakebox}

\textit{Rules callout ``CALCULATE, INTERPRET, THEN BOUND THE CLAIM (keep on the page)'' is printed on the student sheet.}\par\medskip

\dokbadge{1}~\textbf{(a)} Calculate $\hat p_A$, $\hat p_B$, and the observed difference $\hat p_A-\hat p_B$.\par
\answer{The sample proportions are $\hat p_A=3050/5000=0.61$ and $\hat p_B=2950/5000=0.59$, so the observed difference is $0.61-0.59=0.020$, or 2 percentage points.}

\dokbadge{2}~\textbf{(b)} Verify conditions and carry out the two-sided two-sample $z$-test of $H_0:p_A-p_B=0$. Compute $\hat p_c$, pooled SE, $z$, and the $p$-value; interpret the $p$-value and decide at $\alpha=0.05$.\par
\answer{Random selection and assignment support independent groups; $10{,}000\le0.10(500{,}000)=50{,}000$. The pooled proportion is $0.60$, giving expected counts 3,000 and 2,000 per group. The pooled SE is $0.00979796$, $z=2.04124$, and the two-sided $p$-value is $0.04123$. If the population proportions were equal, the chance of a sample difference at least $0.020$ in absolute value is about $0.041$. Since $0.04123<0.05$, reject $H_0$; there is evidence the proportions differ.}

\dokbadge{3}~\textbf{(c)} Decide which report is warranted. Use the $p$-value, observed gap, 5-point standard, and null value. Explain what the rejected report would mislead a reader to believe and name the confidence level and endpoint criterion needed for deployment.\par
\answer{Salma is warranted. The test is significant because $0.04123<0.05$, but the observed difference $0.020$ is below $0.050$, and the test used null value 0. It therefore does not establish an improvement above 5 points. Rowan conflates statistical detectability with practical sufficiency. At a 5 percent one-sided level, deployment would require a 95 percent lower confidence bound for $p_A-p_B$ above $0.05$. Random selection and assignment support population and causal inference, but do not change the claim tested.}

\begin{summaryexitbox}{EXIT}
One thing the video changed about my first take: \hrulefill
\end{summaryexitbox}

\end{document}
